\documentclass[11pt,a4paper]{article}
\usepackage[T1]{fontenc}
\usepackage{lmodern,amsmath,amssymb,amsthm,booktabs,array}
\usepackage[margin=26mm]{geometry}
\usepackage{xcolor}
\usepackage[colorlinks=true,linkcolor=blue!40!black,citecolor=blue!40!black,urlcolor=blue!40!black]{hyperref}
\newtheorem{theorem}{Theorem}
\newtheorem{lemma}[theorem]{Lemma}
\newtheorem{definition}[theorem]{Definition}
\newcommand{\NP}{\mathsf{NP}}
\newcommand{\FP}{\mathsf{FP}}
\newcommand{\UPAC}{\mathsf{UPAC}}
\newcommand{\SAT}{\mathsf{SAT}}
\newcommand{\R}{\mathbb R}
\newcommand{\N}{\mathbb N}
\newcommand{\Z}{\mathbb Z}
\setlength{\emergencystretch}{2em}
\setlength{\parskip}{0.25em}
\title{Unconstrained Detection of Productive\\Autocatalytic Cores Is NP-Complete}
\author{}
\date{18 September 2026}
\begin{document}
\maketitle
\begin{abstract}
We prove NP-completeness of deciding whether a finite reversible reaction
network contains a productive autocatalytic core, when neither a target
species nor an allowed-food set is specified. The input is a binary encoding
of the two literal nonnegative integer complex matrices, and each reaction
identity carries a signed real flow. The proof first characterizes existence
of a core by a nonsingular square restriction with literal two-sided
participation. A weighted incidence construction then reduces prescribed
two-path directed linkage to this unrestricted search: incidence saturation
controls every possible selected support, and a boundary determinant
distinguishes the desired pairing from the crossed pairing. An explicit
constant-size switch construction reduces SAT to the required linkage
instances. Polynomial certificates and a total polynomial-time binary
transducer complete the classification. The terminal theorem, including
actual machine-level membership and reduction bounds, is verified in Lean.
\end{abstract}

\section{The decision problem}
Selecting the internal species is a substantive part of detecting
autocatalysis. A network can fail to amplify one chosen set while containing
a different productive core. We study the existential question over all
nonempty species and reaction supports.

Kosc et al.~\cite{kosc} define productive motifs and their inclusion-minimal
cores. Their complexity theorem imposes a target species and an allowed-food
set; the published article explicitly leaves the unconstrained variant
open. We resolve that question for finite reversible sources with integer
complexes. Andersen et al.~\cite{andersen} consider a different
resource-production formulation of autocatalysis. Recent algorithms for
enumerating autocatalytic subsystems in large networks~\cite{enumeration}
also address practical instances; worst-case decision complexity does not
preclude effective enumeration on particular networks.

\begin{definition}
A \emph{source} consists of finite sets $E$ of entities and $R$ of reaction
identities, and matrices $L,P\in\N^{E\times R}$. Its net matrix is
$N=P-L$, interpreted over $\R$. A candidate is a pair of nonempty subsets
$(X,S)$ of $E$ and $R$. It is \emph{admissible} if, for every $r\in S$,
\begin{equation}\label{eq:literal}
 \exists e\in X\; L_{er}>0,
 \qquad \exists e\in X\; P_{er}>0.
\end{equation}
It is \emph{productive} if there is $v\in\R^R$, zero outside $S$, such that
\begin{equation}\label{eq:production}
 \sum_{r\in R}N_{er}v_r>0\qquad(e\in X).
\end{equation}
A \emph{motif} is an admissible productive candidate. A \emph{productive
autocatalytic core} (PAC) is a motif minimal under componentwise inclusion:
there is no motif $(X',S')$ with $X'\subseteq X$, $S'\subseteq S$, and at
least one inclusion strict.
\end{definition}

Flows may be negative: reversing a reaction changes the sign of its flow,
rather than introducing a second independently selectable identity.
Condition~\eqref{eq:literal} uses the literal matrices even when an entity
occurs on both sides. Cancellation in $N$ must not change admissibility.
Entities outside $X$ have no production constraint. The input contains no
distinguished target and no list of allowed external entities.

We use a canonical dense binary encoding of $|E|$, $|R|$, and the entries
of $L$ and $P$ in reaction-major order. Natural-number fields use binary
digits with delimiters; malformed strings are rejected. Precise conventions
and the formal interfaces appear in Appendix~\ref{sec:formal}.
Let $\UPAC$ be the language of well-formed sources containing a PAC.

\begin{theorem}\label{thm:main}
The language $\UPAC$ is NP-complete under polynomial-time many-one reductions.
\end{theorem}

The theorem concerns structural possibility in this literal-source model.
It does not impose conserved elemental compositions, elementary reaction
arity, concentration bounds, or a kinetic law. In particular, a structural
YES certificate alone does not establish a dynamically realizable pathway.
For network screening, the definitions specify exactly what a certificate
checks: literal participation and strictly positive internal production.

\section{Square witnesses and membership in NP}
Finiteness immediately implies that every motif contains a PAC. The useful
stronger fact is that a PAC has a square nonsingular net restriction.

\begin{lemma}\label{lem:square}
A finite source contains a PAC if and only if it has nonempty subsets
$X\subseteq E$, $S\subseteq R$ such that $|X|=|S|$, every reaction in $S$
satisfies~\eqref{eq:literal}, and $N_{X,S}$ is nonsingular.
\end{lemma}
\begin{proof}
Let $(X,S)$ be a PAC and choose a productive flow $v$. Every $v_r$, $r\in S$,
is nonzero, since otherwise deleting $r$ preserves a motif. Orient each
selected reaction in the direction of its flow, and write $w_r=|v_r|>0$.
The following observation gives $|X|\leq |S|$: each $e\in X$ is the sole
internal reactant of some oriented selected reaction.

Indeed, suppose this fails for $e$. Remove $e$, and delete each reaction
whose only internal product is $e$. Every remaining reaction still has an
internal reactant, by the supposition, and an internal product, by the rule
for deletion. For any remaining entity, a deleted reaction had nonpositive
net contribution. Thus removing those contributions preserves its strictly
positive production. The remaining entity set is nonempty: every original
reaction had a reactant other than $e$. Strict production on this set also
forces the remaining reaction set to be nonempty. We have a proper motif,
a contradiction. Reactions witnessing the observation for distinct entities
are distinct, proving the claimed cardinality inequality.

Next, the columns of $A=N_{X,S}$ are linearly independent. Otherwise choose
$z\ne0$ with $Az=0$, and choose $r$ with $z_r\ne0$. The modified flow
$v-(v_r/z_r)z$ has the same production and is zero at $r$. Deleting $r$
again yields a proper motif; the selected reaction set cannot become empty
because production remains strictly positive. Therefore $|S|\leq |X|$,
and equality follows. The square matrix $A$ is nonsingular.

Conversely, if the stated square restriction exists, solve $Av=\mathbf1$
and extend $v$ by zero outside $S$. This is a motif. A minimal motif below
it exists by finiteness, and is a PAC.
\end{proof}

This argument allows literal catalysts. Only in a source with disjoint
reactant and product supports does admissibility become the assertion that
every selected reaction column of $N$ has both signs.

\pagebreak[3]
\begin{lemma}\label{lem:np}
$\UPAC\in\NP$.
\end{lemma}
\begin{proof}
For a YES instance choose the square restriction in Lemma~\ref{lem:square}.
Let its order be $q$ and let $H\geq1$ bound the absolute values of its integer
entries. With $d=\det A\ne0$, the integer vector
\[
 z=\operatorname{sgn}(d)\operatorname{adj}(A)\mathbf1
 \quad\text{satisfies}\quad Az=|d|\mathbf1>0.
\]
The determinant expansion bounds every cofactor by $(q-1)!H^{q-1}$,
including the usual empty-minor convention when $q=1$. Hence
$|z_r|\leq q!H^{q-1}$, so each signed entry has
$O(q\log(q+1)+q\log(H+1))$ bits. In a nonempty dense input of length $n$,
$q\leq n$ and $\log(H+1)=O(n)$, giving a polynomial-size certificate.

The certificate supplies $X,S$ and this integer flow, extended by zero.
A verifier parses the source, checks both literal sides for every selected
reaction, and performs exact integer matrix--vector multiplication and
strict comparisons on the selected entities. All operand lengths and loop
counts are polynomial in $n$. Accepting such a certificate establishes a
motif and therefore PAC existence; the verifier need not test minimality.
Malformed sources are rejected. The formal proof implements the verifier
as a nondeterministic machine and proves its polynomial time bound.
\end{proof}

\section{From directed linkage to an arbitrary selected support}
Consider a directed graph with distinct terminals $a_0,a_1,b_0,b_1$. We ask
for vertex-disjoint directed paths $a_0\leadsto b_0$ and $a_1\leadsto b_1$.
We may delete arcs entering the sources, arcs leaving the sinks, and loops:
none is needed by a pair of simple disjoint paths. The remaining vertices
are called internal. Parallel arcs, if present, retain distinct identities.

Construct an integer matrix $M$ with one column for each arc and one row
for each internal vertex, together with two boundary rows $s,t$. Both
sources are identified with row $s$, and both sinks with row $t$. For an
arc $e:u\to w$, only its tail row and head row are nonzero. Define
\[
 \alpha(a_0)=-1,\quad\alpha(a_1)=1,\quad\alpha(u)=-1
 \quad(u\text{ internal}),
\]
\[
 \beta(b_0)=-1,\quad\beta(b_1)=1,\quad\beta(w)=1
 \quad(w\text{ internal}).
\]
Let $f(a_1)=f(b_1)=2$, with $f=1$ at all other endpoints. Set
\begin{equation}\label{eq:arc}
 M_{\mathrm{tail}(e),e}=\alpha(u),\qquad
 M_{\mathrm{head}(e),e}=\beta(w)f(u)f(w).
\end{equation}
The two rows are distinct, and every nonzero coefficient has magnitude at
most four. Turn $M$ into a source by taking arcs as entities, matrix rows
as reactions, and
\begin{equation}\label{eq:source}
 L_{er}=\max\{-M_{re},0\},\qquad
 P_{er}=\max\{M_{re},0\}.
\end{equation}
Thus the net matrix is $M^{\mathsf T}$ and the literal sides are disjoint.

\begin{lemma}\label{lem:link}
The source~\eqref{eq:source} contains a PAC if and only if the graph has
vertex-disjoint paths $a_0\leadsto b_0$ and $a_1\leadsto b_1$.
\end{lemma}
\begin{proof}
First consider any square nonsingular witness supplied by
Lemma~\ref{lem:square}, transposed to a $q\times q$ restriction of $M$.
Every selected row is mixed, so it has at least two selected nonzero
entries. Every selected column has at most two nonzero entries. Counting
incidences gives
\begin{equation}\label{eq:saturation}
 2q\ \leq\ \#\{(r,e):M_{re}\ne0\text{ in the restriction}\}\ \leq\ 2q.
\end{equation}
Equality forces both endpoints of every selected arc to be selected, and
exactly two selected arcs, of opposite signs, at every selected row.
At an internal vertex these are one entering and one leaving arc. At $s$
they are one arc from each of $a_0,a_1$; at $t$ they are one arc into each
of $b_0,b_1$.

If neither boundary row is selected, the finite selected graph is a union
of directed cycles on internal vertices. Each such cycle has ordinary
$-1,+1$ incidence coefficients and makes the restriction singular. If $s$
is selected, following either outgoing path through the unique outgoing
arc at every internal vertex cannot enter a cycle: entry would violate
uniqueness of the incoming arc. It must reach a sink. The reverse argument
applies if $t$ is selected. Thus nonsingularity forces both boundary rows,
two vertex-disjoint source-to-sink paths, and possibly separate internal
cycles. Those cycles again give a singular block and are excluded.

There are two possible pairings of the paths. Eliminate the internal rows
along each path, using their nonzero incidence pivots. Apart from nonzero
row or column scaling, the remaining boundary blocks are
\begin{equation}\label{eq:blocks}
 B_{\mathrm{desired}}=
 \begin{pmatrix}-1&1\\-1&4\end{pmatrix},
 \qquad
 B_{\mathrm{crossed}}=
 \begin{pmatrix}-1&1\\2&-2\end{pmatrix}.
\end{equation}
For example, the factor two leaving $a_1$ propagates along its path, and
the factor two entering $b_1$ multiplies it once more. This also covers a
direct source-to-sink arc by~\eqref{eq:arc}. The determinants are $-3$ and
$0$. Consequently a nonsingular selection must use the desired pairing.

Conversely, take a desired pair of vertex-disjoint simple paths. Select
their arcs, their internal vertices, and $s,t$. Each path has one more arc
than internal vertices, so the selection is square. Every selected row is
mixed. Internal elimination leaves $B_{\mathrm{desired}}$, so the selected
matrix is nonsingular. Lemma~\ref{lem:square} supplies a PAC.
\end{proof}

The quantifier over all restrictions is essential here. Equation~\eqref{eq:saturation}
prevents a witness from using only one endpoint of an arc, branching at an
internal vertex, or retaining a boundary-free productive component. No
promise that a candidate includes the full intended gadget is required.

\section{An explicit reduction from SAT to linkage}
We give the graph construction used by the formal reduction, so the
hardness argument does not rely on an assumed linkage-completeness theorem.
Paths in this section are simple. A CNF formula has $c$ clauses, $o$ literal
occurrences, and $V$ variable slots, including any unused indices. Index
occurrences by $0\leq i<o$, retaining repeated literals as separate
occurrences, and put $K=o+1$.

\subsection{The switch}
Use the fixed directed graph on vertices $0,\ldots,19$ listed in
Appendix~\ref{sec:switch}. Its input ports are $0,1,2,3$ and its output
ports are $4,5,6,7$. It has the following three properties.

\begin{lemma}\label{lem:switch}
For the fixed switch:
\begin{enumerate}
\item If disjoint paths connect an input $a$ to $4$ and $1$ to an output
$b$, then $a=0$ and $b=5$.
\item Given disjoint control paths $0\leadsto4$ and $1\leadsto5$, every
input-to-output path disjoint from both controls is either
$2\leadsto6$ or $3\leadsto7$. Both kinds cannot exist for the same controls,
even without requiring the two residual paths to be disjoint from one another.
\item There are disjoint controls with an available $2\leadsto6$ path,
and there are disjoint controls with an available $3\leadsto7$ path.
\end{enumerate}
\end{lemma}
\begin{proof}
Appendix~\ref{sec:switch} specifies the complete finite check of the first
two assertions, together with the completeness argument for its path
enumeration. The third assertion has explicit witnesses in the same
appendix. All three assertions are proved in Lean; the finite checks are
kernel-reduced decisions, not external numerical tests.
\end{proof}

\subsection{Stack and formula wiring}
Take $K$ switches, indexed by $i=0,\ldots,K-1$, writing $(i,p)$ for port
$p$ at level $i$. Join
\begin{equation}\label{eq:stack}
 (i,5)\longrightarrow(i+1,1),\qquad
 (i+1,4)\longrightarrow(i,0)\qquad(0\leq i<K-1).
\end{equation}
The last switch is a dummy level with no literal occurrence. Introduce
variable gates $u_0,\ldots,u_V$, rail vertices $h_{j,b,k}$ for
$0\leq j<V$, $b\in\{0,1\}$, $0\leq k\leq K$, and clause gates
$d_0,\ldots,d_c$. Add the following arcs, and no others beyond switch
and stack arcs.

First, $(K-1,5)\to u_0$, and for each variable $j$ add
$u_j\to h_{j,b,0}$ for both choices of $b$. Between consecutive positions
on a rail, use $h_{j,b,i}\to h_{j,b,i+1}$ unless occurrence $i$ is a literal
of variable $j$ that is false when $j=b$. In that case replace the step by
\begin{equation}\label{eq:rail}
 h_{j,b,i}\longrightarrow(i,3),\qquad
 (i,7)\longrightarrow h_{j,b,i+1}.
\end{equation}
The dummy position always uses the direct rail step. Add
$h_{j,b,K}\to u_{j+1}$, and $u_V\to d_0$. Finally, for each occurrence $i$
in clause $\ell$, add
\begin{equation}\label{eq:clause}
 d_\ell\longrightarrow(i,2),\qquad
 (i,6)\longrightarrow d_{\ell+1}.
\end{equation}
The prescribed linkage is
\begin{equation}\label{eq:terminals}
 (0,1)\leadsto d_c,
 \qquad (K-1,0)\leadsto(0,4).
\end{equation}
The first path is denoted $P$ and the second $Q$. The dummy level makes
$K>0$ for every formula; the four terminals remain distinct in boundary cases.

\begin{lemma}\label{lem:sat}
The graph above has the linkage~\eqref{eq:terminals} if and only if the
CNF formula is satisfiable.
\end{lemma}
\begin{proof}
Suppose first that an assignment satisfies the formula. At a true
occurrence choose switch controls leaving channel $2\leadsto6$ available;
at a false occurrence choose controls leaving $3\leadsto7$ available.
Either choice suffices at the dummy level. The path $Q$ follows the upward
controls $0\leadsto4$, and the initial part of $P$ follows the downward
controls $1\leadsto5$. After the stack, $P$ chooses the assignment's rail
for each variable. Exactly the false occurrences require a detour through
$3\leadsto7$. Then, in each clause, choose a true occurrence and traverse
its available $2\leadsto6$ channel. Distinct occurrences are separate
switches. A true occurrence is not used on a false-literal rail detour,
so the clause and rail portions do not conflict. These choices give the
two disjoint paths.

For the converse, let $P,Q$ be any disjoint paths with the required
endpoints. We first force the controls at every level, without assuming
that the paths respect their intended roles. At level zero, $P$ contains
port $1$ and $Q$ contains port $4$ by their endpoints. Extract the segment
of $P$ from that $1$ to its next exit from the switch, and the segment of
$Q$ from its last entrance into the switch to that $4$. Entrances use input
ports and exits use output ports. The two segments are disjoint, so
Lemma~\ref{lem:switch}(1) forces $P$ to exit at $5$ and $Q$ to enter at $0$.
The unique continuation from $(i,5)$ is $(i+1,1)$, and the unique predecessor
of $(i,0)$ is $(i+1,4)$, except at the stack boundaries.
Induction therefore forces the downward control in $P$ and the upward
control in $Q$ at every level. Simplicity gives the corresponding order
of their visits.

After $P$ leaves $(K-1,5)$, its remaining suffix is disjoint from its own
control prefix and from $Q$. Any later input-to-output passage through a
switch must consequently be a residual channel. By
Lemma~\ref{lem:switch}(2), it preserves the rail or clause role: it cannot
enter at $3$ and leave at $6$, or enter at $2$ and leave at $7$.
The external wiring now forces this suffix to pass the variable gates in
order, choose one rail for each variable, and subsequently pass the clause
gates in order. Its rail choices define an assignment. Its passage between
$d_\ell$ and $d_{\ell+1}$ chooses an occurrence in clause $\ell$ and uses
that occurrence's $2\leadsto6$ channel. If the chosen literal were false
under the assignment, the earlier rail passage would have used the same
switch's $3\leadsto7$ channel. This contradicts residual exclusivity for
the fixed controls. Every chosen clause literal is therefore true.

An empty clause has no connecting occurrence and prevents the clause
passage. With no clauses, the clause portion is empty and the construction
accepts. Thus the argument includes these boundary cases.
\end{proof}

\section{Binary execution and completion of the proof}
For the SAT encoding used by the formalization, occurrence count, clause
count, and the largest variable index plus one are bounded by a polynomial
in the input length $n$ (indeed by linear bounds in the chosen encoding).
The stack and rail construction therefore has polynomially many vertices
and arcs. Equation~\eqref{eq:arc} has bounded coefficients, and the dense
source~\eqref{eq:source} has polynomial encoding length. Its entries can be
generated by bounded loops over the explicit vertex and arc families.

To obtain a language reduction on \emph{all} binary strings, parse the
input CNF. On a valid input, write the source constructed above. On a
malformed input, write the source for the fixed unsatisfiable formula
consisting of one empty clause. Lemmas~\ref{lem:link} and~\ref{lem:sat}
show that this fixed output is a NO instance and that the resulting total
function $F$ satisfies
\begin{equation}\label{eq:reduction}
 x\in\SAT\quad\Longleftrightarrow\quad F(x)\in\UPAC
 \qquad\text{for every binary string }x.
\end{equation}

For completeness, the formal execution proof goes beyond bounding output
size. A parser prepares a validity flag and payload; a selector returns
either the decoded formula's encoding or the fixed unsatisfiable encoding.
The valid-input writer starts on that ordinary input with blank work tapes,
writes the headers, and traverses both literal matrix sides in canonical
order. Its loop invariants and a polynomial bound on its actual transitions
are proved. A composition theorem combines these runs, including copying
and tape-boundary overhead, into a total transducer computing $F$ in $\FP$.
No well-formed-input premise remains in its final interface.

\begin{proof}[Proof of Theorem~\ref{thm:main}]
Lemma~\ref{lem:np} gives membership in $\NP$. The total polynomial-time
map $F$, with equivalence~\eqref{eq:reduction}, reduces SAT to $\UPAC$.
SAT hardness follows from the Cook--Levin theorem; the formal development
proves and imports that theorem rather than postulating SAT hardness.
Polynomial-time reduction closure yields NP-hardness, and hence
NP-completeness.
\end{proof}

\appendix
\section{The fixed switch and its complete check}\label{sec:switch}
This appendix specifies the only finite gadget check used in the proof.
The vertex set is $\{0,\ldots,19\}$ and the complete arc set is
\begin{align*}
 &(0,9),(0,15),(1,17),(1,18),(2,8),(3,12),\\
 &(8,5),(8,9),(9,10),(9,11),(10,8),(10,11),\\
 &(11,6),(11,19),(12,13),(13,5),(13,14),(14,12),\\
 &(14,15),(15,7),(15,16),(16,4),(16,10),(17,4),\\
 &(17,18),(18,16),(18,19),(19,14),(19,17).
\end{align*}
Let $\mathcal E(k,H,a)$ enumerate simple path lists starting at $a$, with
at most $k$ vertices and avoiding the previously visited set $H$. Its
recursive definition is
\[
\mathcal E(k,H,a)=
\begin{cases}
 \varnothing,& k=0\text{ or }a\in H,\\
 \{[a]\}\ \cup\displaystyle\bigcup_{a\to b}
 \{a::p:p\in\mathcal E(k-1,H\cup\{a\},b)\},&\text{otherwise}.
\end{cases}
\]
Every simple path has at most 20 vertices. Induction on its length shows
it belongs to $\mathcal E(20,\varnothing,a)$. Filtering these lists by
last vertex therefore gives all paths between any two specified ports.
The first two assertions of Lemma~\ref{lem:switch} are finite universal
tests over these lists, using disjointness of vertex sets. In particular,
for each disjoint pair of controls, the exclusivity test checks that it is
not the case that a residual $2\leadsto6$ path exists \emph{and} a residual
$3\leadsto7$ path exists. This stronger statement is what the formula
converse consumes.

In Lean, \texttt{enumerate\_complete} proves coverage of the enumeration,
and \texttt{path\_mem\_routes} transfers arbitrary simple paths into it.
The Boolean identities \texttt{controlCheck\_true},
\texttt{residualCheck\_true}, and \texttt{exclusiveCheck\_true} are proved
by \texttt{decide}. Their logical consequences are
\texttt{control\_ports}, \texttt{residual\_channel}, and
\texttt{residual\_exclusive}. Thus the finite reduction is part of the
kernel proof, with no unproved search-coverage assumption.

The constructive switch modes are displayed below. Within each column the
three paths are pairwise vertex-disjoint; every consecutive pair is an
arc in the displayed graph.
\begin{center}
\small
\begin{tabular}{@{}lll@{}}
\toprule
Role & Channel $2\leadsto6$ available & Channel $3\leadsto7$ available\\
\midrule
$0\leadsto4$ & $0,15,16,4$ & $0,9,11,19,17,4$\\
$1\leadsto5$ & $1,17,18,19,14,12,13,5$ & $1,18,16,10,8,5$\\
Residual & $2,8,9,10,11,6$ & $3,12,13,14,15,7$\\
\bottomrule
\end{tabular}
\end{center}

\section{Formal theorem and reproducibility}\label{sec:formal}
The mathematical source files are under
\texttt{proofs/UnconstrainedPACDetection/}. The terminal declaration is
\begin{center}
\small\ttfamily
UnconstrainedPACDetection.complexity\_root\par
\normalfont of type\par
\ttfamily Complexity.NPComplete\par
UnconstrainedPACDetection.BinaryPACVerifier.language
\end{center}
It has no hypotheses. Its reported axioms are exactly
\texttt{propext}, \texttt{Classical.choice}, and \texttt{Quot.sound}.
There are no admitted lemmas or assumed mathematical hardness axioms in
this interface. Strict verification treats warnings as errors.

The source encoding writes the two dimensions followed by all entries of
the left matrix and then all entries of the right matrix, with reaction
index outermost and entity index innermost. A natural-number field uses
its canonical least-significant-bit-first binary digits, each prefixed
by $1$, followed by a terminating $0$. Canonical decoding rejects invalid
representations and incorrect field counts. The formal language accepts
precisely decoded sources with a PAC. In particular, a net matrix is not
silently substituted for the pair of literal matrices.

\begin{center}
\small
\begin{tabular}{@{}p{.35\linewidth}p{.59\linewidth}@{}}
\toprule
Mathematical obligation & Principal formal component\\
\midrule
Core and square equivalence & \texttt{CoreWitness}, \texttt{SquareMinor}\\
Linkage in both directions & \texttt{LinkagePACForward}, \texttt{LinkagePACConverse}\\
SAT linkage equivalence & \texttt{FormulaLinkage.formula\_paths\_iff\_sat}\\
All-string semantic reduction & \texttt{FormulaPACEncoding.compile\_correct}\\
Actual total FP transducer & \texttt{FormulaTotalWriter.compile\_mem\_FP}\\
Actual NP membership & \texttt{VerifierNPPolynomial.in\_NP}\\
Final classification & \texttt{complexity\_root} in \texttt{ComplexityRoot.lean}\\
\bottomrule
\end{tabular}
\end{center}

The frozen toolchain is \texttt{leanprover/lean4:v4.30.0}, with dependencies
fixed by \path{mathlib4_project/lake-manifest.json}. The strict receipt is
\path{complexity_root.verify.json}, stored in the workspace's
\path{campaign_2026_09_05/} directory. It records the source and transitive dependency hashes,
the authenticated declaration types, and the axiom report. Publication
verification rechecks those hashes; it is distinct from claiming a fresh
clean build or a performance benchmark.

From the repository root, the verification bridge can be invoked as follows
(the long declaration is one command argument):
\begin{quote}\small\ttfamily\raggedright
.venv/Scripts/python.exe scripts/verify\_proof.py\par
\quad proofs/UnconstrainedPACDetection/ComplexityRoot.lean\par
\quad --timeout 600\par
\quad --declaration UnconstrainedPACDetection.complexity\_root
\end{quote}
The bridge runs Lean in the pinned project and checks the dependency
closure. The publication directory includes the editable source,
build script, and evidence audit. The complete proof is provided in the
repository accompanying this manuscript; no external archival DOI or
submission is asserted.

\begin{thebibliography}{9}
\bibitem{kosc}
T. Kosc, D. Kuperberg, E. Rajon, and S. Charlat.
Thermodynamic consistency of autocatalytic cycles.
\emph{Proceedings of the National Academy of Sciences},
122(18):e2421274122, 2025.
\href{https://doi.org/10.1073/pnas.2421274122}{doi:10.1073/pnas.2421274122}.

\bibitem{andersen}
J. L. Andersen, C. Flamm, D. Merkle, and P. F. Stadler.
Maximizing output and recognizing autocatalysis in chemical reaction
networks is NP-complete.
\emph{Journal of Systems Chemistry}, 3:1, 2012.
\href{https://doi.org/10.1186/1759-2208-3-1}{doi:10.1186/1759-2208-3-1}.

\bibitem{enumeration}
R. Golnik, T. Gatter, P. F. Stadler, and N. Vassena.
Enumeration of autocatalytic subsystems in large chemical reaction networks.
\emph{Journal of Chemical Theory and Computation}, 22(10):4888--4907, 2026.
\href{https://doi.org/10.1021/acs.jctc.5c01979}{doi:10.1021/acs.jctc.5c01979}.
\end{thebibliography}
\end{document}
